\section{Research problem, objectives and boundaries}
\subsection{Industrial problem}
Conformal coatings are thin protective polymer films placed over electronic assemblies to mitigate moisture, contaminants and related electrical-reliability risks. Their optical transparency makes ordinary visible-light inspection difficult. Where the formulation contains an appropriate fluorescent tracer, illumination in the ultraviolet may reveal coating coverage. However, \emph{visibility of fluorescence is not a direct demonstration of acceptable coating thickness, adhesion, curing state or long-term protection}. IPC-HDBK-830A addresses coating selection, processing, coverage and inspection; IPC-A-610J addresses assembled-board acceptability, including coating coverage and thickness \citep{ipchdbk830a,ipca610j,ipccc830c}. The publicly accessible standard records were reviewed at the level of their tables of contents and summaries. They are not used here to invent numerical acceptance rules.

The target is a nominal 70~$\times$~70~mm populated PCB in repeatable production positioning. A camera viewing an approximately 80~$\times$~80~mm area is a tentative optical design assumption, providing a 5~mm margin around each nominal board edge. A fixed single-camera station may be materially simpler than scanning an entire panel with a moving XY platform. But this does not guarantee detection on component sidewalls, under tall components or behind connectors. Commercial systems explicitly add oblique cameras for those areas \citep{goepelcci,mirteccci}. Accordingly, the research concerns \emph{measurable visible regions under defined optical conditions}, not universal three-dimensional coating verification.

\subsection{Research aim and objectives}
The overarching aim is to assess whether a low-cost, fixed-geometry UV fluorescence station combined with reference-aware image analysis can perform reliable, auditable inspection of selected conformal-coating defects within a production-relevant cycle time. The measurable objectives are to: (i) characterize excitation, emission, camera response and registration error; (ii) define a board-specific coverage map from CAD, the coating plan and expert-labelled visible zones; (iii) compare rule-based and learned segmentation under identical conditions; (iv) identify and size defects with explicit keep-out/required-zone semantics; (v) quantify failure to detect actionable defects; and (vi) measure practical CPU-only end-to-end latency and cost.

A scientifically responsible claim requires a relevant comparator. Prior UV bubble inspection studies already combine classical candidate selection with CNN confirmation \citep{lee2021cnn,lee2021transfer}. Recent studies optimize UV intensity using statistical features and YOLOv8 tiling \citep{choi2024uv}, or train detector networks for dewetting and capillary defects \citep{patel2026coverage}. Therefore \emph{neither hybrid image processing nor CNN use for conformal coatings is automatically novel}. Novelty, if any, must arise from an isolated and measured property such as reference-constrained defect handling, illumination robustness, uncertainty-aware human review, resource efficiency, or validation across manufacturing lots.

\subsection{Research questions and falsifiable hypotheses}
\textbf{RQ1:} Under which controlled lighting and camera settings can UV fluorescence distinguish coating-covered from uncovered \emph{visible} PCB regions with repeatable intensity and acceptable registration error?\par
\textbf{RQ2:} Does a reference-aware hybrid improve \emph{board-level actionable-defect sensitivity at matched false-reject rate} relative to classical thresholding and standalone segmentation?\par
\textbf{RQ3:} Does multi-angle excitation or exposure bracketing reduce misses in shadow-prone regions enough to justify additional exposure/cycle-time cost?\par
\textbf{RQ4:} Can a CPU-class computer support the complete acquisition-to-result workflow within an actual specified production takt requirement, while meeting lighting stability and repeatability requirements?\par
\textbf{RQ5:} How do performance and calibration shift across physical boards, coating batches, operator handling, camera positioning and lighting aging?\par

\textbf{H1 (to test):} a geometric coating plan constrains false alarms from naturally dark components and improves protection of keep-out zones at fixed sensitivity. \textbf{H2 (to test):} flat-field normalization decreases performance variability across illumination geometries. \textbf{H3 (to test):} a lightweight segmentation network with reference-aware rules can be more resource-efficient than a substantially larger end-to-end model. The null hypotheses state that no meaningful improvement exists. Thresholds, comparator configurations and decision criteria must be frozen before evaluating the independent test set.

\subsection{Claims explicitly outside scope}
This manuscript does not assert full-coating thickness metrology by UV brightness; does not equate an image pixel with optical resolving power; does not infer that a 0.1~mm physical void can be detected reliably simply because it spans multiple pixels; and does not estimate a production false-acceptance rate from synthetic images. Destructive or OCT thickness measurements and electrical testing can independently validate selected hypotheses. OCT studies explicitly show how layer interfaces and thickness can be measured when fluorescence alone cannot provide such evidence \citep{shao2019oct,lf2023oct}.
